% ConFu++ paper draft — generated from the repository's real experimental record.
% All numbers trace to EXPERIMENTS.md and results/ JSON artifacts.
\documentclass[10pt,twocolumn]{article}
\usepackage[margin=1in]{geometry}
\usepackage{booktabs,graphicx,amsmath,amssymb,hyperref,multirow}
\newcommand{\rIA}{r_{12}}
\newcommand{\ra}{r_1}
\newcommand{\rb}{r_2}
\title{Alignment Is Not Dependence:\\Measuring and Reducing Unimodal Shortcuts in\\Higher-Order Multimodal Fusion}
\author{Anonymous}
\date{}
\begin{document}
\maketitle

\begin{abstract}
Higher-order multimodal methods fuse representations of several modalities and align them
contrastively, on the premise that fused representations capture information unavailable to any
constituent modality. We show this premise is empirically unfounded. Across three established
benchmarks (AV-MNIST, CMU-MOSI, UR-FUNNY), interaction representations that are active,
predictive, high-rank, and permutation-sensitive are almost entirely predictable from a single
modality (nonlinear predictability up to $R^2=0.989$). We formalize the resulting
\emph{alignment--dependence gap} and introduce an operational diagnostic suite that separates
fusion, dependence, and conditional utility. We then propose \textbf{S1}, a training-time
adversarial de-shortcutting objective that penalizes predictability of each interaction from each
single modality while leaving joint predictability unconstrained. S1 is the only objective among
those studied that significantly reduces the measured shortcut (paired $p\approx0.011$, five
seeds) at zero accuracy cost, and it fabricates no dependence on negative controls. On a new
controlled synthetic benchmark with ground-truth third-order synergy, order-restricted models sit
at chance while an order-matched interaction solves the task with perfectly balanced modality
dependence ($+49.9$~pp, 5/5 seeds). Finally, on a newly released natural benchmark with genuine
headroom (Bird-MML), the original higher-order objective captures none of the available
complementarity ($+0.2$~pp of $+15.9$~pp available), while our utility+S1 stack recovers a
significant share ($+5.8$~pp over baseline, paired $p<0.001$).
We release the diagnostic suite, the synthetic benchmark, and all experimental logs.
\end{abstract}

\section{Introduction}
Multimodal learning increasingly relies on fusing per-modality encodings into joint
representations, aligned with contrastive objectives (CLIP-style and higher-order variants).
The implicit claim: the fused representation \emph{integrates} the modalities. We ask a simpler
question that the literature rarely answers directly --- does the fused representation actually
\emph{depend} on all of its inputs?

Our central finding is a repeatable negative: fused interactions can satisfy every surface health
criterion (variance, effective rank, standalone probes, permutation sensitivity) while being a
unimodal function. We make this precise with predictability probes: on AV-MNIST,
$R^2(\ra\!\to\!\rIA)=0.989$ while $R^2(\rb\!\to\!\rIA)=0.031$.

\textbf{Contributions.}
(1) We formalize the alignment--dependence gap and release an operational diagnostic suite
(\S\ref{sec:diagnostics}). (2) We audit three standard benchmarks and show their higher-order
interactions are shortcut-dominated (\S\ref{sec:audit}). (3) We propose S1, an adversarial
de-shortcutting objective with a hinged, joint-sparing penalty; it is the only tested objective
that significantly reduces single-modality predictability at no accuracy cost, and it is honest
on negative controls (\S\ref{sec:s1}). (4) We contribute a controlled synthetic order benchmark
on which order-matched interactions are provably (empirically) necessary and sufficient, while
order-restricted decompositions sit at chance (\S\ref{sec:orderbench}). (5) We document two
methodology pitfalls: post-aggregation LayerNorm leaks higher-order signal, invalidating
``low-order'' ablations; and online adversaries equilibrate far below converged probes (the
\emph{adversary gap}) (\S\ref{sec:pitfalls}).

\section{Related Work}
% See docs/RELATED_WORK.md for the full annotated survey (verification pass pending).
\paragraph{Contrastive multimodal alignment.} CLIP and its tri-modal extensions (Symile, TriCLIP,
ImageBind) align embeddings pairwise or via total correlation. ConFu (our direct baseline) aligns
fused pair representations with the remaining modality. None of these objectives constrains the
\emph{content} of the fused representation.
\paragraph{Modality dominance.} Multimodal networks overfit (Wang et al.), gradient rebalancing
(OGM-GE and follow-ups), and underspecification (D'Amour et al.) address dominance at the
loss/gradient level. We show why loss-level pressure cannot bootstrap dependence: corrupted and
positive interactions coincide under a shortcut.
\paragraph{Synergy and information decomposition.} PID (Williams--Beer) and MultiViz-style
interaction auditing motivate our vocabulary; our metrics are operational approximations, not
information-theoretic quantities.
\paragraph{Adversarial representation control.} Domain-adversarial training and
``learning-not-to-learn'' remove unwanted predictability. S1 differs in target: predictability of
the \emph{interaction} from \emph{each constituent alone}, never from the pair jointly.

\section{The Alignment--Dependence Gap}\label{sec:diagnostics}
\textbf{Setup.} Encoders $E_i$ produce $r_i=E_i(X_i)$; a fusion module produces
$\rIA=F(\ra,\rb)$. We distinguish four properties: \emph{active} (non-trivial variance),
\emph{predictive} ($I(Y;\rIA)>0$), \emph{dependent} (performance degrades when either constituent
is corrupted), and \emph{conditionally useful}
($\mathrm{Perf}(\ra,\rb,\rIA)>\mathrm{Perf}(\ra,\rb)$). Only the combination of the last two
earns the name \emph{complementary}.

\textbf{Diagnostics.} (i) zero vs.\ shuffle interaction tests; (ii) per-modality corruption drops
$D_1, D_2$ and the balance score $B_{12}=2\min(D_1,D_2)/(D_1+D_2)$; (iii) capacity-matched
conditional linear and nonlinear probes; (iv) correction--regression counts against the
lower-order model; (v) nonlinear predictability probes $q_1(\ra)\!\to\!\rIA$,
$q_2(\rb)\!\to\!\rIA$, $q_{12}(\ra,\rb)\!\to\!\rIA$ with the shortcut criterion
$R^2(q_1)\!\gg\!R^2(q_2)$ and $R^2(q_1)\!\approx\!R^2(q_{12})$.

\section{Auditing Existing Benchmarks}\label{sec:audit}
\begin{table}[h]\centering\small
\caption{Five-seed audits. Interactions are healthy on the surface and shortcut underneath.}
\begin{tabular}{lll}
\toprule
Benchmark & Surface health & True content \\
\midrule
AV-MNIST & rank ${\uparrow}$, shuffle-sensitive & $R^2(\ra{\to}\rIA){=}0.989$, audio drop $+0.004$ \\
CMU-MOSI & text dominance & pair conditional gains $-0.96/-1.43/-2.27$~pp; $R^2\!\le\!0.953$ \\
UR-FUNNY & high rank, sensitive & $R^2\!\approx\!0.97$; pair gains $-0.66/{+}0.19/{-}0.76$~pp \\
\bottomrule
\end{tabular}
\end{table}

Every audited benchmark fails either the dependence gate (a modality shuffle is free) or the
utility gate (lower-order probes already saturate). These are \emph{data properties}; no training
objective can create the missing information, and we log all failed attempts as negative results.

\section{S1: Adversarial De-Shortcutting}\label{sec:s1}
\textbf{Mechanism.} For each interaction $t=\mathrm{std}(\rIA)$ (EMA statistics, stop-grad),
adversaries $q_1, q_2$ are trained to predict $t$ from $\ra$ alone and $\rb$ alone (alternating
minimax, $k_{\mathrm{adv}}{=}1$). The generator pays the hinged penalty
$\mathrm{ReLU}(\tau-\mathrm{MSE}(q_i(r_i), t))$, taxing single-modality $R^2$ above $1{-}\tau$
without demanding noise. Joint predictability from $(\ra,\rb)$ is never penalized.
Per-factor variance floors close the constant-factor escape route.

\textbf{Results (AV-MNIST, five seeds, paired).} Accuracy $70.05{\pm}0.48$ vs.\ $69.96{\pm}0.59$
(ablation, ns); $R^2(\ra{\to}\rIA)$: $0.989\to0.682$ ($t=-4.43$, $p\approx0.011$); audio
dependence doubles ($+0.084\to+0.236$~pp, $p\approx0.059$). S1 is the first objective in this
study to move single-modality predictability at all; utility-loss (P1) and ranking (P2) variants
never did.

\textbf{Honesty.} On synthetic no-synergy controls and on a distractor-modality control, S1
fabricates nothing (drops $\approx0$; gate stays silent). On AV-MNIST the residual
$R^2\approx0.68$ is an adversary gap (\S\ref{sec:pitfalls}), not hidden dependence.

\section{The Controlled Order Benchmark}\label{sec:orderbench}
Three (to six) modalities observe latent bits through fixed random prototypes plus Gaussian noise;
labels are $y=\bigoplus_i x_i$ (order $M$) or $y=x_1{\oplus}x_2$ with distractors (control).
Probes use held-out splits; five seeds per cell.

\begin{table}[h]\centering\small
\caption{Order-3 synergy ($y{=}a{\oplus}b{\oplus}c$): accuracy (5 seeds).}
\begin{tabular}{lcc}
\toprule
Model & orders & accuracy \\
\midrule
Additive & 1 & $0.496\pm0.006$ \\
Pairs-only & 1+2 & $0.501\pm0.010$ \\
Concat-MLP & joint & $1.000\pm0.000$ \\
\textbf{Full (order-matched)} & 1+2+3 & $\mathbf{1.000\pm0.000}$ \\
\bottomrule
\end{tabular}
\end{table}

The $M$-ladder ($M{=}3..6$) confirms the pattern at every $M$ (order-restricted: chance;
order-matched: $1.0000$; concat-MLP cracks at $M{=}6$, $0.9975$). The full model's dependence is
maximally balanced (every single-modality shuffle costs ${\approx}0.5$ accuracy), and
held-out probes of singles/pairs sit at chance. Notably, pair representations do contain the
bits (concat probe 0.97) --- the order-2 failure is a \emph{readout} failure of sum aggregation,
evidence that higher-order synergy requires order-matched interaction \emph{computation}, not
merely rich features.

\section{Bird-MML: A Natural Benchmark with Genuine Headroom}\label{sec:birdmml}
Bird-MML (149 species; image, audio, text) is audited on frozen pretrained features with a fixed
stratified split (71-species subset; the Zenodo release is missing one archive part --- all
reported factors are identical across compared methods). The audit finds exactly one strong pair:
image+text reaches $80.89\%$ vs.\ best single $64.95\%$ ($+15.94$~pp headroom), while audio
carries nothing in speech-domain features.

\begin{table}[h]\centering\small
\caption{Bird-MML, five seeds paired. ConFu++ = utility objective + S1.}
\begin{tabular}{lccc}
\toprule
 & ConFu (orig.) & ConFu+U & ConFu+U+S1 \\
\midrule
probe full acc (\%) & $76.90\pm0.65$ & $83.01\pm0.13$ & $82.66\pm0.37$ \\
conditional gain (pp) & $+0.19$ & $+1.07$ & $+0.77$ \\
$R^2(\mathrm{text}{\to}z_{13})$ & $0.59$ & $\mathbf{0.44}$ & $0.50$ \\
audio shuffle drop (pp) & $+0.01$ & $+0.01$ & $+0.01$ \\
\bottomrule
\end{tabular}
\end{table}

The original objective leaves the available complementarity unused (full ${\approx}$ lower).
ConFu++ recovers $+5.76{\pm}0.86$~pp over it (paired $t{=}14.91$, $p{<}0.001$), exceeding the
raw-feature linear ceiling ($81.01\%$), while fabricating no dependence for the null modality.

\paragraph{Fair-play controls.} (i) \textbf{Leakage control:} $33\%$ of released captions
contain the species name; re-auditing with BLIP-only (name-free) captions reduces the headroom
from $+15.9$ to $+4.8$~pp --- still positive, so genuine complementarity exists beyond leakage.
(ii) \textbf{Attribution:} the accuracy gain comes from the utility objective (ConFu+U
$83.01{\pm}0.13$ vs.\ +S1-only $76.58{\pm}0.51$); S1's role is representation quality, not
score --- on the balanced pair, text-predictability of $z_{13}$ is lowest under the utility
objective ($R^2$ $0.59{\to}0.44$) and S1 pins it to the hinge band ($0.50$ at $\tau{=}0.5$).
(iii) \textbf{Same-topic baseline:} OGM-style gradient rebalancing on identical frozen features
does not beat plain joint training ($80.68{\pm}0.38$ vs.\ $80.88{\pm}0.11$, paired
$t{=}{-}1.16$, ns) and sits $1.8$~pp below ConFu++. (iv) \textbf{Hyperparameter robustness:}
S1 is insensitive on ground-truth synergy (synthetic Config A: acc ${\approx}1.0$, balanced
drops, for $\tau{\in}\{0.3,0.5,0.7\}$, $k_{adv}{\in}\{1,3\}$).

\section{End-to-End Replication}\label{sec:e2e}
To verify the frozen-feature findings transfer, we train encoders, fusion, and objectives jointly
on the same subset (identical initialization, schedule, and inference capacity across arms; three
seeds). Every arm converges to \textbf{text dominance} (text shuffle drop ${+}4.9$~pp, image
${\approx}0$, audio $0$), because the released captions contain species names (33\%) and the
task loss absorbs the leakage. Accuracy is statistically flat across arms (confu $71.87{\pm}0.16$,
+U $72.18{\pm}0.48$, +U+S1 $72.32{\pm}0.62$; paired $+0.45{\pm}0.49$~pp, ns), while the
shortcut meters respond monotonically as designed: $R^2(\mathrm{text}{\to}z_{13})$ falls
$0.81 \to 0.73 \to 0.68$ across confu $\to$ +U $\to$ +U+S1. End-to-end, the
alignment--dependence gap manifests as \emph{modality abandonment} --- a stronger failure mode
than the frozen-feature shortcut, and exactly what our diagnostics are built to expose.

\section{Methodology Pitfalls}\label{sec:pitfalls}
\textbf{Post-aggregation LayerNorm is a hidden higher-order term.} A final LayerNorm divides by a
cross-modal statistic; additive models ending in LayerNorm reach $0.75$ on 3-bit parity (exactly
$0.50$ without). ``Low-order'' ablations must use LN-free readouts.
\textbf{The adversary gap.} An online adversary updated once per batch equilibrates at
$R^2\approx0.3$--$0.4$ (below our hinge band) while a converged offline probe reads $0.84$;
de-shortcutting claims must therefore be measured with converged probes, not training-time meters.

\section{Discussion}
The causal chain is now complete: where higher-order information is absent or dominated
(AV-MNIST, MOSI, UR-FUNNY), interactions shortcut; where it is present and observable (synthetic
order benchmark), order-matched interactions extract it completely and S1 neither breaks nor
fabricates anything. The immediate future work is a natural benchmark with validated headroom
(Bird-MML), where the same audit-then-de-shortcut pipeline applies unchanged.
\textbf{Limitations.} S1's hinge tolerates redundancy by design ($\tau$); pair-conditional
utility of individual interaction terms remains negative (utility arrives via the joint head);
the Bird-MML subset uses frozen features (an end-to-end replication trains overnight at writing)
and the audio channel is null under speech-domain features (an AudioSet-domain re-extraction is
in progress); caption name leakage is documented ($33\%$) with a leakage-free control arm.
All predictability statements are operational approximations.

\section*{Reproducibility}
Five seeds, paired statistics, capacity-matched probes, train-only parameter accounting,
git-state logging per run, and public release of the diagnostic suite, synthetic benchmark,
configs, and experiment logs.

\bibliographystyle{plain}
\begin{thebibliography}{9}\small
\bibitem{confu} Koutoupis et al. The More, the Merrier: Contrastive Fusion for Higher-Order
Multimodal Alignment. CVPR 2026 / arXiv:2511.21331.
\bibitem{clip} Radford et al. Learning transferable visual models from natural language
supervision. ICML 2021.
\bibitem{symile} Meyer et al. Symile: A total correlation objective for multimodal contrastive
learning. ICML 2024.
\bibitem{imagebind} Girdhar et al. ImageBind: One embedding space to bind them all. CVPR 2023.
\bibitem{multibench} Liang et al. MultiBench: Multiscale benchmarks for multimodal representation
learning. NeurIPS 2021.
\bibitem{multiviz} Liang et al. MultiViz: An analysis benchmark for multimodal interactions.
2022.
\bibitem{pid} Williams \& Beer. Nonnegative decomposition of multivariate information. 2010.
\bibitem{wang2020} Wang et al. What makes training multi-modal classification networks hard?
CVPR 2020.
\bibitem{ogm} Peng et al. Balanced multimodal learning via on-the-fly gradient modulation. CVPR
2022.
\bibitem{underspec} D'Amour et al. Underspecification presents challenges for credibility in
modern machine learning. JMLR 2022.
\bibitem{dann} Ganin \& Lempitsky. Unsupervised domain adaptation by backpropagation. ICML 2015.
\bibitem{learningnottoplearn} Kim et al. Learning not to learn: Training deep neural networks with
biased data. CVPR 2019.
\bibitem{vicreg} Bardes et al. VICReg: Variance-invariance-covariance regularization. ICLR 2022.
\bibitem{lmf} Liu et al. Efficient low-rank multimodal fusion. ACL 2018.
\bibitem{focal} Lin et al. Focal loss for dense object detection. ICCV 2017.
\end{thebibliography}
\end{document}
